\section{Cross-Compilation Example}

\textbf{Same TypeScript code, three architectures:}

\begin{lstlisting}[style=ts]
import {
  LLVMCodegen,
  ARMCodegen,
  WASMCodegen,
  LLVMFunction,
  BasicBlock,
  RetInst
} from "@euriklis/llvm-ir";

function createAddFunction(codegen) {
  const func = new LLVMFunction({
    name: "add",
    returnType: LLVMCodegen.types.i32
  });
  // ... (same function body for all architectures)
  codegen.getModule().addFunction(func);
}

// x86-64
const x86 = new LLVMCodegen("x86");
createAddFunction(x86);
console.log("=== x86-64 ===");
console.log(x86.toString());

// ARM64
const arm = new ARMCodegen("arm", true);
createAddFunction(arm);
console.log("=== ARM64 ===");
console.log(arm.toString());

// WebAssembly
const wasm = new WASMCodegen("wasm");
createAddFunction(wasm);
console.log("=== WASM ===");
console.log(wasm.toString());
\end{lstlisting}

Each will have a different target triple, but the function \acr{IR} is identical. \acr{LLVM}'s backends handle architecture-specific optimizations (register allocation, instruction selection, calling conventions).

